\bookchapter{Capítulo 29: Salud digital — Estableciendo una relación saludable con el móvil}

\chapterintro{¿Te ha pasado?: solo querías ver vídeos cortos 5 minutos y, al levantar la vista, han pasado dos horas. O estás en la cama por la noche con el teléfono y cada vez estás más despierto, hasta que son las 2 de la madrugada. El teléfono ha cambiado nuestra vida, pero también se está tragando en silencio nuestro tiempo y nuestra salud. Este capítulo no trata de “dejar el teléfono”, sino de recuperar el control sobre él.}

\booksection{I. Mira tu “tiempo de pantalla” - los números te abrirán los ojos}

Antes de cambiar, entiende la situación actual. El teléfono tiene una función que mucha gente nunca ha abierto, llamada “Tiempo de uso” (iPhone) o “Bienestar digital” (Android), que te dice cuánto tiempo pasas en el teléfono cada día y en qué apps.

\begin{itemize}
  \item \textbf{Sinceramente, es normal que el móvil empiece a irse lento tras un ano de uso. La causa no necesariamente es que este roto; más a menudo es "acumulacion de basura": archivos basura, procesos en segúndo plano, espacio de almacenamiento llenandose poco a poco.}: Ajustes > Tiempo de uso > Ver toda la actividad. Verás el número de desbloqueos diarios, cuántas veces coges el teléfono y el tiempo de uso de cada app.
  \item \textbf{Android}: Ajustes > Bienestar digital y control parental. Huawei lo llama "uso saludable del teléfono"; Xiaomi "gestión de tiempo de pantalla". Funciones similares, todas muestran datos de uso diario.
\end{itemize}

\textbf{El 80\% de quienes ven estos datos por primera vez se asustan}. Un amigo mío descubrió que pasaba más de 4 horas diarias en Douyin, 28 horas a la semana, el equivalente a un trabajo a tiempo parcial. Este “impacto digital” es el primer paso para el cambio.

\booksection{II. Establece “límites de uso” para las apps - ponte un limitador de velocidad}

Una vez sepas dónde está el problema, puedes establecer límites específicos. No necesitas borrar las apps; solo ponles un “límite diario”:

\begin{enumerate}
  \item \textbf{Sinceramente, es normal que el móvil empiece a irse lento tras un ano de uso. La causa no necesariamente es que este roto; más a menudo es "acumulacion de basura": archivos basura, procesos en segúndo plano, espacio de almacenamiento llenandose poco a poco.}: Ajustes > Tiempo de uso > Límite de apps > Añadir límite. Selecciona la app a limitar (como Douyin) y establece 30 minutos diarios. Al llegar al límite, el icono se vuelve gris y abrirlo requiere una contraseña adicional.
  \item \textbf{Android}: Ajustes > Bienestar digital > Temporizador de apps. Selecciona una app y establece un límite diario; al llegar al límite, el icono se vuelve gris.
  \item \textbf{Truco}: Pon una contraseña de límite que no uses habitualmente y que no sea fácil de recordar. Si es demasiado fácil de introducir, siempre "ampliarás" el límite y será inútil.
\end{enumerate}

\booksection{III. El uso correcto del “modo no molestar” - no es apagar notificaciones, es apagar la ansiedad}

A mucha gente le suena el teléfono decenas de veces al día con notificaciones; cada vez, la atención se interrumpe. Crees que solo “miras una notificación”, pero en realidad, volver a concentrarte después de cada interrupción requiere 15-20 minutos. El modo no molestar no consiste en “perder el contacto”, sino en recuperar el control de tu atención.

\begin{itemize}
  \item \textbf{Activación programada}: Configura el modo no molestar para que se active automáticamente de 10 de la noche a 7 de la mañana. Así el teléfono no te despertará con un SMS basura mientras duermes.
  \item \textbf{Permitir llamadas importantes}: En los ajustes del modo no molestar, pon "permitir llamadas" como "contactos favoritos" o "llamadas repetidas" (una segunda llamada en 3 minutos sonará normalmente). Así no te perderás una llamada de emergencia de la familia, pero tampoco te molestarán los comerciales.
  \item \textbf{Modo de concentración de trabajo}: El "Modo de concentración" del iPhone permite crear escenarios personalizados; por ejemplo, el modo "trabajo" solo permite notificaciones de apps relacionadas con el trabajo, y el modo "lectura" silencia todas las notificaciones.
\end{itemize}

\booksection{IV. Cuidar los ojos no depende de apps sino de hábitos - la regla 20-20-20}

Mirar la pantalla del teléfono durante mucho tiempo causa sequedad y cansancio ocular inevitables. El modo de protección ocular (filtro de luz azul) ayuda algo, pero lo realmente efectivo es cambiar el hábito visual:

\begin{tcbinfo}{Regla 20-20-20}
Cada\textbf{20 minutos}mirando la pantalla, desvía la mirada a algo a\textbf{20 pies (unos 6 metros)}de distancia y mantenla durante\textbf{20 segundos}. Este simple gesto alivia eficazmente la fatiga del músculo ciliar.

\end{tcbinfo}

\begin{itemize}
  \item \textbf{Activar el modo de protección ocular}: Ajustes > Pantalla > Modo de protección ocular (o modo nocturno); al activarlo, la temperatura de color se vuelve más cálida, reduciendo la luz azul. Recomendado activarlo automáticamente de sunset a sunrise.
  \item \textbf{Ajustar el brillo de la pantalla}: Si el entorno es oscuro, la pantalla también debe estar oscura. Mucha gente pone el brillo demasiado alto en la oscuridad; esto es lo más dañino para los ojos. Activa el "brillo automático" y asegúrate de que funciona.
  \item \textbf{Mantén la distancia}: El teléfono debe estar al menos a 30 cm de los ojos. A mucha gente le gusta sostener el teléfono muy cerca; así el músculo ciliar está en constante tensión y la miopía aumenta rápidamente.
\end{itemize}

\booksection{V. Guía de salvamento cervical para “adictos a mirar hacia abajo”}

Al mirar el teléfono hacia abajo, tu cuello soporta 5-6 veces la presión de estar de pie. Mantener esta postura a largo plazo causa rectificación de la curvatura cervical, proliferación ósea, mareos y entumecimiento de manos. Aquí hay algunos métodos sencillos de mejora:

\begin{itemize}
  \item \textbf{Levanta el teléfono}: Al mirar el teléfono, intenta levantarlo a la altura de los ojos, en lugar de bajar la cabeza para mirar el teléfono sobre las piernas. Aunque los brazos se cansen, tu cuello te lo agradecerá.
  \item \textbf{Mueve el cuello cada 30 minutos}: Gira lentamente la cabeza hacia la izquierda hasta el límite, mantén 3 segundos, vuelve al centro; luego gira a la derecha, mantén 3 segundos. Repite 5 veces. Luego haz el "ejercicio del carácter 米": escribe el carácter 米 en el aire con la barbilla.
  \item \textbf{De pie contra la pared 3 minutos}: Apoya la cabeza, las escápulas, los glúteos y los talones contra la pared; permanece de pie 3 minutos al día para corregir eficazmente la postura encorvada.
\end{itemize}

\booksection{VI. Sueño y teléfono - el hábito que más deberías cambiar en el dormitorio}

La luz azul de la pantalla del teléfono inhibe la secreción de melatonina, haciendo que “cuanto más miras, más despierto estás”. Mucha gente cree que “lo voy mirando y me da sueño”, pero en realidad la luz azul engaña al cerebro y, cuando dejas el teléfono, te cuesta más quedarte dormido. Recomendaciones de mejora:

\begin{itemize}
  \item \textbf{No mires el teléfono durante la hora previa a dormir}: Lo más efectivo pero difícil de cumplir. Si no puedes, al menos activa el "modo nocturno/protección ocular" para reducir la luz azul al mínimo.
  \item \textbf{Pon el teléfono fuera del alcance de la mano}: Coloca el cargador en el otro extremo del dormitorio, no en la mesita de noche. Así no tendrás la tentación de cogerlo y, cuando suene la alarma, tendrás que levantarte para apagarla; matas dos pájaros de un tiro.
  \item \textbf{Activa el "modo sueño"}: El "Modo sueño" del iPhone puede bloquear automáticamente la pantalla durante las horas de descanso y simplificar la pantalla de bloqueo; las marcas Android tienen funciones similares.
  \item \textbf{Usa un despertador en lugar del teléfono}: Compra un despertador físico barato; si no usas el teléfono como despertador, te librarás de la trampa de "mirar la hora y luego no poder evitar mirar el teléfono".
\end{itemize}

\booksection{VII. Ansiedad informativa - aprende a tener “ceguera selectiva”}

La cantidad de información que recibimos cada día es la que una persona antigua recibía en toda su vida. Pero el cerebro humano no está diseñado para procesar tanta información; la sobrecarga causa ansiedad, dispersión de la atención y fatiga de decisión. La solución no es “obtener más información”, sino “bloquear activamente”:

\begin{itemize}
  \item \textbf{Dejar de seguir}: Revisa todas las cuentas públicas, UP masters y bloggers que sigues; pregúntate: "¿He obtenido algún valor real de ellos en el último mes?" Si no, deja de seguir. Mejor pocas fuentes pero de calidad.
  \item \textbf{Desactivar notificaciones de apps no esenciales}: Excepto WeChat, teléfono, SMS y calendario, desactiva las notificaciones de todas las demás apps. Las noticias no necesitan notificaciones; cuando vayas a buscarlas, las verás.
  \item \textbf{Establece un "día de ayuno informativo"}: Elige un día a la semana (como el domingo) sin mirar noticias del teléfono ni vídeos cortos. Descubrirás que el cielo no se cae y el lunes estarás con mejor estado mental.
\end{itemize}

\begin{tcbpractice}{Práctica: reto de bienestar digital de una semana}
\begin{enumerate}
  \item Hoy: abre Tiempo de uso / Bienestar digital y revisa tus datos de uso de la semana pasada.
  \item Mañana: establece un límite diario para la app de entretenimiento que más tiempo consume (recomendado 30-60 minutos).
  \item Pasado mañana: programa el modo no molestar (22:00-07:00) y permite llamadas de contactos favoritos.
  \item Cuarto día: no mires el teléfono durante la hora previa a dormir; pon el cargador en el otro extremo del dormitorio.
  \item Quinto día: deja de seguir 5 cuentas públicas o bloggers que ya no miras.
  \item Sexto día: intenta un "ayuno informativo" de un día entero; sin noticias ni vídeos cortos.
  \item Séptimo día: revisa el informe de tiempo de pantalla de esta semana y compáralo con la semana anterior.
\end{enumerate}

\end{tcbpractice}

\begin{tcbsummary}{Resumen del capítulo}
\begin{itemize}
  \item Primero mira tus datos de tiempo de pantalla; reconocer la situación actual es el primer paso para el cambio
  \item Pon límites a apps de entretenimiento y programa el modo no molestar; recupera el control de tu atención
  \item La regla 20-20-20 es el hábito de cuidado ocular más sencillo y efectivo
  \item Mirar el teléfono hacia abajo multiplica por 5-6 la carga del cuello; levanta el teléfono y mueve el cuello periódicamente
  \item No mires el teléfono durante la hora previa a dormir; pon el cargador fuera del alcance de la mano
  \item La solución para la ansiedad informativa es "bloquear activamente", no "obtener más"
\end{itemize}

\end{tcbsummary}
